\documentclass[../../main.tex]{subfiles}
\begin{document}

\begin{flushright}
\footnotesize\textit{【自動文字起こし・要確認】}
\end{flushright}

[解] (1) 題意のようになるのは、1人目がどちらの席にすわるか場合分けして考えて、確率 $P_1$ として
\[
P_1 = \dfrac{1}{12} \cdot \dfrac{2}{10} + \dfrac{2}{12} \cdot \dfrac{1}{8} = \dfrac{9}{240} = \dfrac{3}{80} \quad \text{（終）}
\]

\paragraph{(2)}

席のうまる順番は
\begin{align}
& 1^\circ \ (2, 4, 6) \quad (6, 4, 2) \cdots \text{\textcircled{1}} \\
& 2^\circ \ (2, 6, 4) \quad (6, 2, 4) \cdots \text{\textcircled{2}} \\
& 3^\circ \ (4, 2, 6) \quad (4, 6, 2) \cdots \text{\textcircled{3}}
\end{align}
2と6の対称性から、\textcircled{1}, \textcircled{2}, \textcircled{3} の2つの順番がおこる確率は、いずれも等しく、排反だから、これらのおこる確率の和がもとめる確率 $P_2$ である\\
\textcircled{1} の時 $\dfrac{2}{12} \cdot \dfrac{2}{9} \cdot \dfrac{2}{6} = \dfrac{1}{81}$ \\
\textcircled{2} の時 $\dfrac{2}{12} \cdot \dfrac{2}{9} \cdot \dfrac{2}{6} = \dfrac{1}{81}$ \\
\textcircled{3} の時 $\dfrac{2}{12} \cdot \dfrac{2}{8} \cdot \dfrac{2}{6} = \dfrac{1}{72}$ \\
より、
\[
P_2 = 2 \left( \dfrac{1}{81} + \dfrac{1}{81} + \dfrac{1}{72} \right) = \dfrac{25}{324} \quad \text{（終）}
\]

\end{document}
